% Fichier généré automatiquement par tools/generate_corrections.py.
% Source : PERF/PERF-05-Precistion-TVF/501_Divers/501_Divers.tex
% Statut : complet (2/2 question(s) renseignée(s), 0 reprise(s)).
\begin{corrigebox}[Corrigé extrait de la banque]
\CorrigeQuestion{1}
On a 
$H(p)=\dfrac{\dfrac{K}{p\left(1+\tau_1 p \right)\left(1+\tau_2 p \right)}}{1+\dfrac{CK}{p\left(1+\tau_1 p \right)\left(1+\tau_2 p \right)}}$
$=\dfrac{K}{p\left(1+\tau_1 p \right)\left(1+\tau_2 p \right)+CK}$. 
En conséquence, $S(p)=E(p)\dfrac{K}{p\left(1+\tau_1 p \right)\left(1+\tau_2 p \right)+CK}$.

$s_{\infty}=\lim\limits_{t\to +\infty} s(t)$ $=\lim\limits_{p\to 0} pS(p)$
$=\lim\limits_{p\to 0} pE(p)H(p)$.
Dans le cas où $E(p)$ est un échelon, on a $E(p)=\dfrac{E_0}{p}$ et donc 
$s_{\infty}=\lim\limits_{p\to 0} p\dfrac{E_0}{p}\dfrac{K}{p\left(1+\tau_1 p \right)\left(1+\tau_2 p \right)+CK}=\dfrac{E_0}{C}$.
\CorrigeQuestion{2}
On a maintenant $E(p)=\dfrac{k}{p^2}$. 
On a donc et donc 
$s_{\infty}=\lim\limits_{p\to 0} p\dfrac{k}{p^2}\dfrac{K}{p\left(1+\tau_1 p \right)\left(1+\tau_2 p \right)+CK}$ et 
$s_{\infty}=\infty$.

\begin{enumerate}
    \item $s_{\infty}=\lim\limits_{p\to 0} p\dfrac{E_0}{p}\dfrac{K}{p\left(1+\tau_1 p \right)\left(1+\tau_2 p \right)+CK}=\dfrac{E_0}{C}$.
   
    \item $s_{\infty}=\infty$.
   

\end{enumerate}
\end{corrigebox}
